\documentclass[11pt]{article}
\usepackage[spanish,es-nodecimaldot,es-noshorthands]{babel}
\usepackage[utf8]{inputenc}
\usepackage[T1]{fontenc}
\usepackage{booktabs}
\usepackage{graphicx}
\usepackage{amsmath}
\usepackage{pgfplots}
\usepackage{placeins}
\usepackage[hidelinks]{hyperref}
\usepackage[a4paper,margin=2.5cm]{geometry}
\pgfplotsset{compat=1.18}
\setlength{\emergencystretch}{2em}
\renewcommand{\arraystretch}{1.12}
\title{Lab 01: Métodos no supervisados}
\author{Dylan Elizondo Alvarado (504610652)\\
Joseph Garcia Montero (402680547)\\
Rodrigo Ureña Castillo (118910482)\\
Mariana Ureña Padilla (119180324)}
\date{EIF420O -- 2026}
\begin{document}
\maketitle

\section{Introducción y objetivos}
Se analiza el conjunto de datos proporcionado para el laboratorio,
conservado en \texttt{data/ejemplo\_analisis.csv}. Contiene 16 observaciones,
cuatro variables numéricas (\texttt{edad}, \texttt{ingresos},
\texttt{visitas} y \texttt{satisfaccion}) y la categoría \texttt{segmento}.
No se dispone de un diccionario que establezca moneda, período de visitas,
escala de satisfacción o población de origen; por ello, se interpretan
relaciones dentro de la muestra, sin atribuirles un contexto comercial no
documentado.

El objetivo es caracterizar los datos mediante análisis exploratorio (EDA),
comparar ACP, HAC, K-Means, t-SNE y UMAP, y seleccionar configuraciones
justificadas por métricas y visualizaciones. La reducción dimensional y el
agrupamiento responden preguntas distintas: representar información en pocos
ejes no equivale a demostrar la existencia de grupos naturales.

\subsection{Exploración de los datos}
Las variables numéricas varían en escalas muy distintas, especialmente
ingresos frente a satisfacción. Hay un faltante en ingresos (fila 6) y otro
en satisfacción (fila 15), con numeración desde 1: dos celdas, en dos filas
distintas. Cada variable afectada tiene 6.25\% de valores ausentes; no se
eliminó ninguna observación. No hay filas completas duplicadas.
La categoría segmento contiene Nuevo (6), Frecuente (6) y Ocasional (4).
Se conserva para describir la entrada, pero se excluye de los modelos y no
se utiliza como verdad de referencia.

\begin{table}[htbp]
\centering\small
\caption{Descriptivos originales, antes de imputar. DE: desviación estándar muestral.}\label{tab:eda}
\begin{tabular}{lrrrrrr}\toprule
Variable & Válidos & Media & Mediana & DE & Mín. & Máx. \\ \midrule
edad & 16 & 38.31 & 38.00 & 11.31 & 22.0 & 56.0 \\
ingresos & 15 & 1876.67 & 1800.00 & 798.95 & 850.0 & 3010.0 \\
visitas & 16 & 8.75 & 8.00 & 5.30 & 2.0 & 18.0 \\
satisfaccion & 15 & 7.82 & 7.70 & 1.01 & 6.2 & 9.3 \\
\bottomrule\end{tabular}
\end{table}
\begin{figure}[htbp]\centering
\begin{tikzpicture}
\begin{axis}[width=7.4cm,height=5.4cm,title={Edad y visitas},xlabel={edad},ylabel={visitas},grid=major,legend style={font=\scriptsize,at={(0.5,-0.24)},anchor=north,legend columns=2},tick label style={font=\scriptsize},title style={font=\small},label style={font=\small}]
\addplot[only marks,mark=*,color=blue!70!black,mark size=2pt] coordinates {(22.000000,2.000000) (24.000000,3.000000) (25.000000,2.000000) (27.000000,4.000000) (29.000000,5.000000) (31.000000,4.000000) (45.000000,12.000000) (47.000000,14.000000) (49.000000,13.000000) (52.000000,16.000000) (54.000000,15.000000) (56.000000,18.000000) (35.000000,7.000000) (37.000000,8.000000) (39.000000,9.000000) (41.000000,8.000000)};
\end{axis}\end{tikzpicture}\hfill
\begin{tikzpicture}
\begin{axis}[width=7.4cm,height=5.4cm,title={Edad e ingresos observados},xlabel={edad},ylabel={ingresos},grid=major,legend style={font=\scriptsize,at={(0.5,-0.24)},anchor=north,legend columns=2},tick label style={font=\scriptsize},title style={font=\small},label style={font=\small}]
\addplot[only marks,mark=*,color=blue!70!black,mark size=2pt] coordinates {(22.000000,850.000000) (24.000000,910.000000) (25.000000,880.000000) (27.000000,1020.000000) (29.000000,1100.000000) (45.000000,2450.000000) (47.000000,2580.000000) (49.000000,2630.000000) (52.000000,2790.000000) (54.000000,2880.000000) (56.000000,3010.000000) (35.000000,1620.000000) (37.000000,1710.000000) (39.000000,1800.000000) (41.000000,1920.000000)};
\end{axis}\end{tikzpicture}
\caption{Relaciones originales. El segundo panel omite la fila cuyo ingreso está ausente.}\label{fig:eda}
\end{figure}

Las seis correlaciones de Pearson, calculadas antes de imputar y con pares
completos, están entre 0.9801 y 0.9942. La relación edad--ingresos es 0.9942;
edad--visitas, 0.9842; edad--satisfacción, 0.9852;
ingresos--visitas, 0.9882; ingresos--satisfacción, 0.9801;
visitas--satisfacción, 0.9873. Son asociaciones descriptivas fuertes, no
evidencia de causalidad. Se anticipa una dirección común dominante para ACP.
Con la regla univariada de 1.5 veces el rango intercuartílico no se detectan
valores atípicos numéricos; esto no descarta anomalías multivariadas.
\FloatBarrier
\clearpage

\section{Fundamentos teóricos}
\subsection{Análisis exploratorio y preparación}
EDA examina distribución, escala, faltantes y relaciones, combinando
estadísticos y gráficos antes de imponer un modelo \cite{nisteda}.
Permite justificar la preparación y detectar problemas que alterarían las
distancias. La estandarización $z_{ij}=(x_{ij}-\mu_j)/\sigma_j$ evita que una
variable domine solamente por sus unidades, aunque no elimina correlaciones
ni vuelve equivalentes sus significados.

\subsection{Análisis de componentes principales (ACP)}
ACP representa los datos centrados mediante ejes ortogonales ordenados por
varianza. Para una matriz preparada $X$, las puntuaciones son $T=XV$ y la
varianza explicada por el componente $j$ es
$\lambda_j/\sum_\ell\lambda_\ell$. Conservar pocos componentes comprime
información, pero puede perder señales de baja varianza \cite{jolliffe2002,jolliffecadima}.
La implementación usa PCA de scikit-learn \cite{sklearnpca}.
El framework exporta cargas $L=V\,\mathrm{diag}(\sqrt{\lambda_j})$:
no son los vectores propios sin escalar ni correlaciones exactas.
El signo de un componente puede invertirse sin cambiar su representación.

\subsection{K-Means}
K-Means alterna asignación al centro más próximo y actualización de
centroides para reducir
\[
  J=\sum_{g=1}^{k}\sum_{i\in C_g}\lVert x_i-\mu_g\rVert^2.
\]
Su inercia mide compactación en el espacio utilizado. Requiere fijar $k$,
puede alcanzar mínimos locales y es sensible a escala y valores extremos
\cite{macqueen1967,sklearnkmeans}. Un $k$ mayor suele reducir la inercia:
su mínimo por sí solo no determina una partición útil.

\subsection{Agrupamiento jerárquico aglomerativo (HAC)}
HAC comienza con observaciones separadas y fusiona grupos hasta construir
una jerarquía \cite{murtagh2012}. Los enlaces single, complete y average
utilizan, respectivamente, la distancia mínima, máxima y promedio entre
pares de puntos de grupos distintos; Ward minimiza el incremento de
variabilidad intragrupo y se utiliza con distancia euclidiana
\cite{scipylinkage}. El dendrograma registra las alturas de fusión; un corte
produce una partición. Las fusiones no se deshacen, por lo que el enlace
puede afectar el resultado.

\subsection{t-SNE}
t-SNE compara probabilidades de vecindad del espacio original y de la
proyección; en el espacio reducido usa una distribución t de colas pesadas
y minimiza una divergencia de Kullback--Leibler \cite{vandermaaten2008}.
Perplexity controla el alcance efectivo de las vecindades y debe ser menor
que el número de observaciones \cite{sklearntsne}. Es un método exploratorio
no lineal: los ejes no representan variables originales y las separaciones
visuales entre grupos no deben leerse como distancias globales preservadas.

\subsection{UMAP}
UMAP construye un grafo ponderado de vecindad y optimiza su representación
en baja dimensión \cite{mcinnes2018}. El número de vecinos regula el énfasis
local frente a uno más amplio; min\_dist controla cuán juntos pueden quedar
los puntos en la representación \cite{umapparams}. Modificar estos parámetros
puede cambiar la apariencia de los grupos, sin demostrar categorías reales.
Se estudia como proyección, no como algoritmo de agrupamiento.

\subsection{Criterios de evaluación}
Silhouette contrasta distancia media dentro del grupo ($a$) y distancia
media al grupo alternativo más cercano ($b$):
$s=(b-a)/\max(a,b)$. Valores altos favorecen separación y cohesión; la media
requiere entre 2 y $n-1$ grupos \cite{sklearnsilhouette}.
Se calcula en las cuatro variables preparadas, no en las proyecciones.

Trustworthiness $T_q$, entre 0 y 1, penaliza vecinos que aparecen en la
proyección sin ser vecinos adecuados en el espacio original
\cite{sklearntrust}. Aquí se fija $q=5$ para comparar proyecciones de las
mismas 16 filas. Un valor alto indica preservación local, no garantiza
conservación global, estabilidad o validez semántica de los grupos.
Silhouette y trustworthiness no son métricas intercambiables.

\begin{center}
\fbox{\parbox{0.90\linewidth}{\centering
\textbf{Mapa conceptual:} datos $\rightarrow$ EDA $\rightarrow$
imputación y escala $\rightarrow$
\{ACP, K-Means, HAC, t-SNE, UMAP\} $\rightarrow$
varianza / silhouette / trustworthiness $\rightarrow$
interpretación, comparación y selección.}}
\end{center}
\clearpage

\section{Metodología reproducible}
\subsection{Integración y ambiente}
Los experimentos se ejecutan con el paquete local \texttt{framework\_ia}.
La clase \texttt{Cluster}, en
\texttt{framework\_ia/modelos/no\_supervisado/agrupamiento.py}, implementa
los cinco métodos y comparte la matriz preparada mediante
\texttt{PreprocesadorNoSupervisado}. El ejecutor
\texttt{scripts/ejecutar\_lab01.py} organiza los barridos y exporta evidencia.

La ejecución registrada usa Python 3.12.14, NumPy 2.5.3, pandas 3.0.6,
SciPy 1.18.1, scikit-learn 1.9.1, umap-learn 0.5.12 y Numba 0.68.0.
Se fija semilla 42 donde el método utiliza aleatoriedad; UMAP usa un
trabajador. Las versiones y el hash SHA-256 de la entrada quedan en
\texttt{resultados/configuracion.json}. Una semilla facilita repetir el
experimento en ese ambiente, pero no prueba estabilidad ante otras
semillas o versiones.

\subsection{Preprocesamiento}
Se seleccionan las cuatro variables numéricas y se mantienen las 16 filas.
Los faltantes se imputan con la \textbf{mediana}: ingresos = 1800 y
satisfacción = 7.7. Después se aplica StandardScaler, calculando medias y
desviaciones poblacionales sobre los datos imputados.
No se codifica segmento, no se descartan columnas constantes porque no las
hay y la matriz final tiene dimensión $16\times4$. El EDA utiliza datos
originales; los modelos y métricas usan la matriz imputada y estandarizada.
Esta distinción evita confundir promedios originales con centroides preparados.

\subsection{Configuraciones y selección}
La configuración base se refiere a la interfaz del framework utilizado,
no a todos los valores predeterminados de las bibliotecas subyacentes.
\begin{table}[htbp]\centering\small
\caption{Configuraciones base y variaciones evaluadas.}
\begin{tabular}{lp{5cm}p{6cm}}\toprule
Método & Base & Variaciones \\ \midrule
ACP & 2 componentes & 2, 3 y 4 componentes \\
K-Means & $k=3$, 20 inicializaciones, 500 iteraciones máximas &
$k=2,\ldots,10$; demás parámetros fijos \\
HAC & Ward, $k=3$, distancia euclidiana &
Ward, average, complete y single; $k=2,\ldots,10$ \\
t-SNE & Perplexity 30, 1000 iteraciones; base operativa 5 &
5 y 15; 30 se omite por $30\not<16$ \\
UMAP & 15 vecinos, min\_dist 0.1, dos dimensiones &
Vecinos $\{5,10,15\}$ y min\_dist $\{0.0,0.1,0.5\}$ \\
\bottomrule\end{tabular}\end{table}

ACP selecciona el menor número \emph{evaluado} de componentes que alcanza
95\% de varianza. K-Means y HAC seleccionan el mayor silhouette; en empate
se prefiere menor $k$ y, para HAC, orden alfabético del enlace.
t-SNE y UMAP seleccionan el mayor $T_5$; se desempata por menor perplexity,
o por menor número de vecinos y después menor min\_dist. Estas reglas
documentadas evitan escoger una figura solamente por su apariencia.

\subsection{Evidencia conservada}
Se generaron 41 CSV y dos JSON en \texttt{resultados}: EDA, matriz preparada,
comparaciones, cargas y varianza ACP, proyecciones, asignaciones, centroides,
perfiles y matrices de vinculación HAC. El resumen registra
\texttt{estado: completado} y la única omisión es t-SNE con perplexity 30.
Los CSV completos permiten verificar valores no redondeados. Las tablas y
figuras de este documento incorporan una instantánea de esa ejecución;
si cambia la entrada o la configuración, deben actualizarse también aquí.
Las instrucciones operativas de instalación y ejecución están separadas
en \texttt{COMO\_EJECUTAR.md}.
\FloatBarrier
\clearpage

\section{Resultados y análisis}
\subsection{ACP}
\begin{table}[htbp]
\centering\small
\caption{Varianza explicada del ACP sobre cuatro variables preparadas.}\label{tab:acp}
\begin{tabular}{lrr}\toprule
Componente & Varianza (\%) & Acumulada (\%) \\ \midrule
CP1 & 98.2950 & 98.2950 \\
CP2 & 1.1479 & 99.4429 \\
CP3 & 0.3439 & 99.7868 \\
CP4 & 0.2132 & 100.0000 \\
\bottomrule\end{tabular}
\end{table}
\begin{table}[htbp]
\centering\small
\caption{Cargas exportadas por el framework para ACP de dos componentes.}\label{tab:cargas}
\begin{tabular}{lrr}\toprule
Variable & Carga CP1 & Carga CP2 \\ \midrule
edad & 1.0291 & 0.0257 \\
ingresos & 1.0179 & 0.1702 \\
visitas & 1.0250 & -0.0851 \\
satisfaccion & 1.0237 & -0.1099 \\
\bottomrule\end{tabular}
\end{table}
\begin{figure}[htbp]\centering
\begin{tikzpicture}
\begin{axis}[width=7.4cm,height=5.4cm,title={ACP: dos componentes},xlabel={CP1},ylabel={CP2},grid=major,legend style={font=\scriptsize,at={(0.5,-0.24)},anchor=north,legend columns=2},tick label style={font=\scriptsize},title style={font=\small},label style={font=\small}]
\addplot[only marks,mark=*,color=blue!70!black,mark size=2pt] coordinates {(-2.943800,0.132780) (-2.395665,-0.176025) (-2.626711,0.037120) (-1.927813,-0.264357) (-1.578954,-0.341383) (-1.278721,0.633586) (-0.656352,0.001864) (-0.300845,-0.064872) (-0.051624,-0.026009) (0.181829,0.035318)};
\addlegendentry{Grupo 0}
\addplot[only marks,mark=triangle*,color=orange!85!black,mark size=2pt] coordinates {(1.320704,0.112020) (1.906667,-0.094148) (1.827860,0.159144) (2.630389,-0.133327) (2.631322,0.108327) (3.261714,-0.120037)};
\addlegendentry{Grupo 1}
\end{axis}\end{tikzpicture}\hfill
\begin{tikzpicture}\begin{axis}[width=7.4cm,height=5.4cm,title={Varianza por componente},xlabel={Componente},ylabel={Varianza (\%)},grid=major,tick label style={font=\scriptsize},title style={font=\small},label style={font=\small}]
\addplot[blue!70!black,mark=*] coordinates {(1.000000,98.294975) (2.000000,1.147884) (3.000000,0.343938) (4.000000,0.213203)};
\end{axis}\end{tikzpicture}
\caption{Proyección ACP y varianza. Los colores identifican los grupos K-Means, no las categorías de segmento.}\label{fig:acp}
\end{figure}

CP1 concentra 98.2950\% de la varianza y tiene cargas positivas similares
en las cuatro variables: representa la dirección compartida de aumento de
edad, ingresos, visitas y satisfacción en esta muestra. CP2 explica apenas
1.1479\%, con una carga positiva mayor en ingresos y negativas en visitas
y satisfacción; distingue variaciones relativas que no siguen exactamente
esa dirección común.
Las cargas pueden superar 1 porque no son correlaciones: StandardScaler
usa desviación poblacional y la varianza del PCA utiliza el divisor $n-1$.

Se seleccionan \textbf{dos componentes}, que conservan 99.4429\%.
Agregar CP3 gana solo 0.3439 puntos porcentuales y CP4 otros 0.2132.
La selección equilibra compresión e interpretación bidimensional. CP1 por
sí solo supera 95\%, pero una sola dimensión no fue parte del barrido y no
se presenta como la configuración seleccionada.
La observación 6 sobresale en CP2 (aproximadamente 0.6336): su ingreso
imputado puede influir en esa desviación, por lo que no se clasifica como
anomalía real sin revisar el dato faltante.
Los colores muestran asignaciones K-Means calculadas en cuatro dimensiones;
ACP no generó esas etiquetas.
\FloatBarrier
\clearpage

\subsection{K-Means}
\begin{table}[htbp]
\centering\small
\caption{Comparación de K-Means. Configuración base: $k=3$.}\label{tab:kmeans}
\begin{tabular}{rrr}\toprule
$k$ & Inercia & Silhouette \\ \midrule
2 & 14.8194 & 0.6261 \\
3 & 5.8810 & 0.5828 \\
4 & 3.8418 & 0.5334 \\
5 & 2.1546 & 0.5029 \\
6 & 1.4148 & 0.5046 \\
7 & 1.1046 & 0.4270 \\
8 & 0.8126 & 0.3819 \\
9 & 0.6090 & 0.3434 \\
10 & 0.4129 & 0.2827 \\
\bottomrule\end{tabular}
\end{table}
\begin{figure}[htbp]\centering
\begin{tikzpicture}\begin{axis}[width=7.4cm,height=5.4cm,title={Compactación},xlabel={$k$},ylabel={Inercia},grid=major,tick label style={font=\scriptsize},title style={font=\small},label style={font=\small}]
\addplot[blue!70!black,mark=*] coordinates {(2.000000,14.819415) (3.000000,5.881012) (4.000000,3.841822) (5.000000,2.154557) (6.000000,1.414805) (7.000000,1.104617) (8.000000,0.812576) (9.000000,0.609035) (10.000000,0.412947)};
\end{axis}\end{tikzpicture}\hfill
\begin{tikzpicture}\begin{axis}[width=7.4cm,height=5.4cm,title={Separación relativa},xlabel={$k$},ylabel={Silhouette},grid=major,tick label style={font=\scriptsize},title style={font=\small},label style={font=\small}]
\addplot[blue!70!black,mark=*] coordinates {(2.000000,0.626073) (3.000000,0.582774) (4.000000,0.533370) (5.000000,0.502946) (6.000000,0.504631) (7.000000,0.427043) (8.000000,0.381928) (9.000000,0.343411) (10.000000,0.282712)};
\end{axis}\end{tikzpicture}
\caption{Curvas de K-Means calculadas en la matriz de cuatro variables estandarizadas.}\label{fig:kmeans}
\end{figure}
\begin{table}[htbp]
\centering\small
\caption{Perfiles medios en unidades originales para K-Means seleccionado.}\label{tab:perfiles}
\begin{tabular}{rrrrrr}\toprule
Grupo & Filas & Edad & Ingresos & Visitas & Satisfacción \\ \midrule
0 & 10 & 31.00 & 1312.22 & 5.20 & 7.12 \\
1 & 6 & 50.50 & 2723.33 & 14.67 & 8.87 \\
\bottomrule\end{tabular}
\end{table}

El máximo silhouette corresponde a \textbf{$k=2$}, con 0.6261.
La base $k=3$ obtiene 0.5828 y reduce la inercia de 14.8194 a 5.8810
(aproximadamente 60.3\%). El codo hace razonable considerar tres grupos
para una descripción más detallada, pero no supera al modelo de dos grupos
en el criterio de separación elegido. Para $k=10$ la inercia llega a
0.4129 y silhouette baja a 0.2827: compactar fragmentando la muestra no
implica mejorar la separación.

Los grupos seleccionados contienen 10 y 6 observaciones. El grupo 1
presenta mayores promedios en todas las variables originales que el grupo 0.
Son perfiles descriptivos de esta entrada; no se les atribuyen categorías
comerciales nuevas. Los perfiles originales omiten cada faltante al calcular
su media: en el grupo 0 los promedios de ingresos y satisfacción utilizan
9 valores, aunque su tamaño es 10. Los centroides preparados, en cambio,
usan todas las filas después de imputar y no están expresados en las unidades
de esta tabla.

No se forzó $k=3$ para coincidir con las tres categorías de segmento:
esa columna no intervino en el ajuste ni en la selección. Los resultados
son internos a la muestra y no constituyen validación externa.
\FloatBarrier
\clearpage

\subsection{HAC}
\begin{table}[htbp]
\centering\small
\caption{Silhouette para las 36 configuraciones HAC. En todas, $k_{\mathrm{real}}=k$.}\label{tab:hac}
\begin{tabular}{rrrrr}\toprule
$k$ & Ward & Average & Complete & Single \\ \midrule
2 & 0.6261 & 0.6261 & 0.6261 & 0.6261 \\
3 & 0.5744 & 0.5828 & 0.5744 & 0.5828 \\
4 & 0.5447 & 0.4924 & 0.5447 & 0.4924 \\
5 & 0.5029 & 0.4714 & 0.5029 & 0.4714 \\
6 & 0.5046 & 0.5046 & 0.5046 & 0.4225 \\
7 & 0.4558 & 0.4558 & 0.4039 & 0.3753 \\
8 & 0.3572 & 0.3640 & 0.3572 & 0.4086 \\
9 & 0.3180 & 0.3278 & 0.2973 & 0.3278 \\
10 & 0.2581 & 0.2671 & 0.2581 & 0.2671 \\
\bottomrule\end{tabular}
\end{table}
\begin{figure}[htbp]\centering
\begin{tikzpicture}\begin{axis}[width=15cm,height=5.5cm,xmin=0.5,xmax=16.5,ymin=0,xlabel={Observación (fila original, numerada desde 1)},ylabel={Distancia de fusión},xtick={1,...,16},xticklabels={7,8,9,12,10,11,4,5,1,2,3,6,16,13,14,15},grid=major,tick label style={font=\scriptsize}]
\addplot[blue!70!black,no marks] coordinates {(15,0.000000) (15,0.293046) (16,0.293046) (16,0.000000)};
\addplot[blue!70!black,no marks] coordinates {(5,0.000000) (5,0.311733) (6,0.311733) (6,0.000000)};
\addplot[blue!70!black,no marks] coordinates {(2,0.000000) (2,0.347939) (3,0.347939) (3,0.000000)};
\addplot[blue!70!black,no marks] coordinates {(7,0.000000) (7,0.357827) (8,0.357827) (8,0.000000)};
\addplot[blue!70!black,no marks] coordinates {(10,0.000000) (10,0.386875) (11,0.386875) (11,0.000000)};
\addplot[blue!70!black,no marks] coordinates {(14,0.000000) (14,0.492759) (15.5,0.492759) (15.5,0.293046)};
\addplot[blue!70!black,no marks] coordinates {(9,0.000000) (9,0.559294) (10.5,0.559294) (10.5,0.386875)};
\addplot[blue!70!black,no marks] coordinates {(1,0.000000) (1,0.576882) (2.5,0.576882) (2.5,0.347939)};
\addplot[blue!70!black,no marks] coordinates {(13,0.000000) (13,0.630237) (14.75,0.630237) (14.75,0.492759)};
\addplot[blue!70!black,no marks] coordinates {(4,0.000000) (4,0.679159) (5.5,0.679159) (5.5,0.311733)};
\addplot[blue!70!black,no marks] coordinates {(7.5,0.357827) (7.5,0.983176) (9.75,0.983176) (9.75,0.559294)};
\addplot[blue!70!black,no marks] coordinates {(1.75,0.576882) (1.75,1.187553) (4.75,1.187553) (4.75,0.679159)};
\addplot[blue!70!black,no marks] coordinates {(12,0.000000) (12,1.289526) (13.875,1.289526) (13.875,0.630237)};
\addplot[blue!70!black,no marks] coordinates {(8.625,0.983176) (8.625,1.963654) (12.9375,1.963654) (12.9375,1.289526)};
\addplot[blue!70!black,no marks] coordinates {(3.25,1.187553) (3.25,3.639737) (10.78125,3.639737) (10.78125,1.963654)};
\addplot[orange!80!black,dashed,no marks] coordinates {(0.5,2.5) (16.5,2.5)};
\end{axis}\end{tikzpicture}
\caption{Dendrograma HAC average, derivado de hac\_linkage.csv. Un corte ilustrativo a 2.5 produce los dos grupos seleccionados.}\label{fig:hac}
\end{figure}

Los cuatro enlaces alcanzan el mismo máximo silhouette, 0.6261, con
\textbf{dos grupos}. Se conserva \textbf{average, $k=2$} por la regla de
desempate alfabética del ejecutor; no existe evidencia métrica de que average
sea superior a Ward, complete o single en ese corte. La base Ward con
tres grupos obtiene 0.5744. Para tres grupos, average y single mejoran a
0.5828, pero siguen debajo del máximo de dos grupos.
En las 36 configuraciones el número efectivo de grupos coincide con el
solicitado; se registra este dato porque un corte maxclust puede dar menos
grupos cuando existen empates de alturas.

En el dendrograma average la penúltima fusión ocurre a 1.9637 y la última
a 3.6397. Un corte entre esas alturas separa un conjunto de 6 observaciones
(filas 7--12) y otro de 10 (filas 1--6 y 13--16). La partición coincide
con K-Means seleccionado, aunque sus identificadores numéricos están
intercambiados. Por tanto, los perfiles son los mismos y no es necesario
reinterpretarlos como grupos diferentes. Esta concordancia entre dos
procedimientos apoya la descripción dentro de la muestra, pero no demuestra
estabilidad ni la existencia universal de dos clases.

Las alturas del dendrograma dependen del enlace y la escala preparada;
no deben compararse como si fueran una misma magnitud entre Ward y average.
Las diferencias para cortes más finos muestran que el enlace sí importa
aunque todos coincidan en el corte de dos grupos.
\FloatBarrier
\clearpage

\subsection{t-SNE}
\begin{table}[htbp]\centering\small
\caption{Comparación t-SNE con cinco vecinos para trustworthiness.}
\begin{tabular}{rrl}\toprule
Perplexity & $T_5$ & Estado \\ \midrule
5 & 0.992188 & Evaluada y seleccionada \\
15 & 0.484375 & Evaluada \\
30 & --- & Omitida: requiere menos de 16 \\
\bottomrule\end{tabular}\end{table}
\begin{figure}[htbp]\centering
\begin{tikzpicture}
\begin{axis}[width=7.4cm,height=5.4cm,title={Perplexity 5},xlabel={TSNE1},ylabel={TSNE2},grid=major,legend style={font=\scriptsize,at={(0.5,-0.24)},anchor=north,legend columns=2},tick label style={font=\scriptsize},title style={font=\small},label style={font=\small}]
\addplot[only marks,mark=*,color=blue!70!black,mark size=2pt] coordinates {(17.717276,16.352869) (13.654519,15.647109) (16.468690,12.864045) (11.880482,12.130684) (8.954843,13.609158) (5.147303,10.489951) (0.734826,9.752356) (-1.503248,6.884070) (-3.793907,9.840516) (-5.558128,6.569969)};
\addlegendentry{Grupo 0}
\addplot[only marks,mark=triangle*,color=orange!85!black,mark size=2pt] coordinates {(-18.399223,1.906184) (-22.448467,2.017766) (-20.670805,-1.080053) (-24.813858,-2.771761) (-26.519115,0.672535) (-28.777363,-2.684981)};
\addlegendentry{Grupo 1}
\end{axis}\end{tikzpicture}\hfill
\begin{tikzpicture}
\begin{axis}[width=7.4cm,height=5.4cm,title={Perplexity 15},xlabel={TSNE1},ylabel={TSNE2},grid=major,legend style={font=\scriptsize,at={(0.5,-0.24)},anchor=north,legend columns=2},tick label style={font=\scriptsize},title style={font=\small},label style={font=\small}]
\addplot[only marks,mark=*,color=blue!70!black,mark size=2pt] coordinates {(26.931580,-15.734031) (-55.313606,157.990890) (84.776020,-76.989914) (36.128986,242.841170) (-50.892540,-25.089685) (-146.132060,89.453255) (132.278170,128.022870) (182.272810,46.095615) (169.149370,-51.171524) (93.772690,33.815900)};
\addlegendentry{Grupo 0}
\addplot[only marks,mark=triangle*,color=orange!85!black,mark size=2pt] coordinates {(18.706718,77.433980) (33.059574,-148.425200) (-141.035550,-32.495335) (-62.595770,59.343760) (-59.092590,-114.353240) (47.158295,157.147280)};
\addlegendentry{Grupo 1}
\end{axis}\end{tikzpicture}
\caption{Comparación t-SNE. Mismos colores K-Means en ambas proyecciones; las escalas de los paneles son independientes.}\label{fig:tsne}
\end{figure}

Se selecciona \textbf{perplexity 5}: preserva vecindades con $T_5=0.9922$,
frente a 0.4844 con 15. La segunda configuración considera una vecindad
muy amplia para 16 filas y, en esta ejecución, conserva peor los vecinos
locales evaluados. El resultado no permite afirmar que perplexity 15 sea
inadecuada para cualquier conjunto de datos.

La base 30 no puede ejecutarse válidamente y se documenta su omisión,
en lugar de sustituirla silenciosamente. La base operativa 5 mantiene el
análisis factible. Los colores K-Means sirven como referencia visual:
no se vuelven a ajustar grupos sobre el plano t-SNE.
La diferencia entre proyecciones respalda la necesidad de comparar
hiperparámetros; separaciones o tamaños visuales no son medidas fiables
de diferencias globales entre grupos. Se mantuvo una sola semilla y no se
comprobó estabilidad al repetir el ajuste con otras.
\FloatBarrier
\clearpage

\subsection{UMAP}
\begin{table}[htbp]
\centering\small
\caption{Trustworthiness de UMAP ($q=5$). $d$ representa min\_dist.}\label{tab:umap}
\begin{tabular}{rrrr}\toprule
Vecinos & $d=0.0$ & $d=0.1$ & $d=0.5$ \\ \midrule
5 & 0.984375 & 0.987500 & 0.984375 \\
10 & 0.987500 & 0.987500 & 0.984375 \\
15 & 0.979688 & 0.967187 & 0.964063 \\
\bottomrule\end{tabular}
\end{table}
\begin{figure}[htbp]\centering
\begin{tikzpicture}
\begin{axis}[width=7.4cm,height=5.4cm,title={UMAP: vecinos 5, distancia 0.1},xlabel={UMAP1},ylabel={UMAP2},grid=major,legend style={font=\scriptsize,at={(0.5,-0.24)},anchor=north,legend columns=2},tick label style={font=\scriptsize},title style={font=\small},label style={font=\small}]
\addplot[only marks,mark=*,color=blue!70!black,mark size=2pt] coordinates {(20.891651,16.333286) (20.477423,16.228674) (20.540997,16.689293) (19.883053,16.145176) (19.629248,15.637339) (19.400105,14.996707) (19.267610,14.348664) (18.969133,13.827809) (19.424652,13.642221) (19.048458,13.299770)};
\addlegendentry{Grupo 0}
\addplot[only marks,mark=triangle*,color=orange!85!black,mark size=2pt] coordinates {(6.381728,-2.367699) (6.767979,-2.972258) (6.862833,-2.579961) (7.542193,-2.984613) (7.458534,-2.397370) (7.971693,-2.740674)};
\addlegendentry{Grupo 1}
\end{axis}\end{tikzpicture}\hfill
\begin{tikzpicture}
\begin{axis}[width=7.4cm,height=5.4cm,title={UMAP: vecinos 15, distancia 0.1},xlabel={UMAP1},ylabel={UMAP2},grid=major,legend style={font=\scriptsize,at={(0.5,-0.24)},anchor=north,legend columns=2},tick label style={font=\scriptsize},title style={font=\small},label style={font=\small}]
\addplot[only marks,mark=*,color=blue!70!black,mark size=2pt] coordinates {(-6.844625,-3.161221) (-6.454635,-3.678986) (-7.056901,-3.588043) (-6.179354,-3.017403) (-6.653932,-2.599771) (-7.369225,-2.342783) (-7.811707,-1.998959) (-8.108183,-3.045131) (-8.227003,-2.573164) (-8.861485,-2.656192)};
\addlegendentry{Grupo 0}
\addplot[only marks,mark=triangle*,color=orange!85!black,mark size=2pt] coordinates {(-9.157450,-1.862733) (-9.565391,-1.399811) (-8.999330,-1.087817) (-9.474153,-0.727848) (-10.041197,-0.876443) (-10.146044,-1.378657)};
\addlegendentry{Grupo 1}
\end{axis}\end{tikzpicture}
\caption{UMAP seleccionado frente a la configuración base. Colores K-Means, escalas independientes.}\label{fig:umap}
\end{figure}

Se selecciona \textbf{5 vecinos y min\_dist 0.1}, con $T_5=0.9875$.
Hay empate con 10 vecinos y min\_dist 0.0 o 0.1; se eligen menos vecinos
según la regla documentada. Para 5 vecinos, min\_dist 0.0 y 0.5 obtienen
0.9844. La base 15/0.1 obtiene 0.9672; reducir vecinos mejora la
preservación local evaluada en esta muestra.

La diferencia entre min\_dist 0.0 y 0.1 no tiene una dirección universal:
con 15 vecinos, 0.0 obtiene 0.9797, mientras que con 5 vecinos gana 0.1.
Esto muestra interacción entre parámetros y evita tratar un valor como
óptimo independiente del contexto. Con 15 vecinos se conecta una fracción
muy grande de las 16 filas; el criterio $T_5$ evalúa una vecindad más local.
Por ello, su selección no equivale a optimizar toda la geometría global.

El t-SNE seleccionado obtiene un $T_5$ ligeramente mayor (0.9922 frente a
0.9875), bajo la misma matriz y vecindad de evaluación. Es una diferencia
descriptiva, no una prueba estadística de superioridad. No se evaluaron
tiempos ni capacidad de generalización. Ninguno de los dos métodos
reemplaza la explicación de las variables que ofrece ACP.
\FloatBarrier

\section{Conclusiones}
El EDA identifica escalas heterogéneas, dos faltantes y fuerte asociación
positiva entre las cuatro variables numéricas. La imputación por mediana y
la estandarización permiten conservar 16 observaciones y comparar métodos
sobre una misma matriz. La categoría segmento queda fuera del ajuste para
mantener el carácter no supervisado.

ACP con dos componentes retiene 99.4429\% y describe una dirección dominante
común. K-Means con dos grupos y HAC average con dos grupos alcanzan
silhouette 0.6261 y la misma partición; el empate de los cuatro enlaces
HAC se resuelve por una convención reproducible, no por superioridad
estadística. La alternativa de tres grupos aporta más detalle y menor
inercia, pero pierde en el criterio elegido. t-SNE con perplexity 5 y UMAP
con 5 vecinos/min\_dist 0.1 preservan bien los vecinos locales de esta
ejecución.

Las conclusiones se restringen a una muestra pequeña, con unidades y
población de origen no documentadas y dos imputaciones. Las métricas son
internas y se evaluó una sola semilla; no hay prueba de estabilidad,
validación externa ni inferencia causal. Un trabajo posterior podría
comparar otras semillas y tratamientos de faltantes, y contrastar los
grupos con información externa antes de utilizarlos para decisiones.
La entrega conserva entrada, código y resultados completos, de modo que
las selecciones pueden verificarse y repetirse.
\clearpage

% Referencias integradas para compilar este documento sin archivos auxiliares.
% referencias.bib se conserva como apoyo para edición bibliográfica en Overleaf.
\begin{thebibliography}{99}
\bibitem{nisteda}
NIST/SEMATECH. \emph{e-Handbook of Statistical Methods:
Exploratory Data Analysis}. \url{https://www.itl.nist.gov/div898/handbook/eda/eda.htm}.

\bibitem{jolliffe2002}
I. T. Jolliffe. \emph{Principal Component Analysis}, segunda edición.
Springer, 2002.

\bibitem{jolliffecadima}
I. T. Jolliffe y J. Cadima. Principal component analysis: a review and
recent developments. \emph{Philosophical Transactions of the Royal Society A},
374:20150202, 2016. \url{https://doi.org/10.1098/rsta.2015.0202}.

\bibitem{macqueen1967}
J. B. MacQueen. Some methods for classification and analysis of multivariate
observations. \emph{Proceedings of the Fifth Berkeley Symposium on
Mathematical Statistics and Probability}, vol. 1, pp. 281--297, 1967.
\url{https://digicoll.lib.berkeley.edu/record/113015}.

\bibitem{murtagh2012}
F. Murtagh y P. Contreras. Algorithms for hierarchical clustering: an
overview. \emph{WIREs Data Mining and Knowledge Discovery}, 2(1):86--97,
2012. \url{https://doi.org/10.1002/widm.53}.

\bibitem{vandermaaten2008}
L. van der Maaten y G. Hinton. Visualizing data using t-SNE.
\emph{Journal of Machine Learning Research}, 9:2579--2605, 2008.
\url{https://www.jmlr.org/papers/v9/vandermaaten08a.html}.

\bibitem{mcinnes2018}
L. McInnes, J. Healy y J. Melville. UMAP: Uniform Manifold Approximation
and Projection for Dimension Reduction, 2018.
\url{https://arxiv.org/abs/1802.03426}.

\bibitem{scipylinkage}
SciPy. \emph{scipy.cluster.hierarchy.linkage}, documentación de referencia.
\url{https://docs.scipy.org/doc/scipy/reference/generated/scipy.cluster.hierarchy.linkage.html}.

\bibitem{sklearnpca}
scikit-learn. \emph{PCA}, documentación de referencia.
\url{https://scikit-learn.org/stable/modules/generated/sklearn.decomposition.PCA.html}.

\bibitem{sklearnkmeans}
scikit-learn. \emph{KMeans}, documentación de referencia.
\url{https://scikit-learn.org/stable/modules/generated/sklearn.cluster.KMeans.html}.

\bibitem{sklearntsne}
scikit-learn. \emph{TSNE}, documentación de referencia.
\url{https://scikit-learn.org/stable/modules/generated/sklearn.manifold.TSNE.html}.

\bibitem{sklearnsilhouette}
scikit-learn. \emph{silhouette\_score}, documentación de referencia.
\url{https://scikit-learn.org/stable/modules/generated/sklearn.metrics.silhouette_score.html}.

\bibitem{sklearntrust}
scikit-learn. \emph{trustworthiness}, documentación de referencia.
\url{https://scikit-learn.org/stable/modules/generated/sklearn.manifold.trustworthiness.html}.

\bibitem{umapparams}
UMAP. \emph{Basic UMAP Parameters}, documentación oficial.
\url{https://umap-learn.readthedocs.io/en/latest/parameters.html}.

\end{thebibliography}
\end{document}
